% =====================================================================
%  Lo que el modelo no representa
%  Anexos
%
%  ENCARGO: Las magnitudes que el modelo deja fuera, cuanto pesan y por
%  que quedan fuera. Nace de la energia reactiva, que ocupaba el 39 % del
%  capitulo 2 para concluir que no altera la comparacion.
%
%  Reglas del documento:
%   - Una idea = una subseccion = una figura.
%   - Toda figura declara su procedencia con.
%   - M1 y M3 se reportan siempre en paralelo.
%   - Ningun superlativo sin releer el CSV de la frontera correcta.
% =====================================================================
\section{Lo que el modelo no representa}
\label{sec:no-representado}
\justifying

El modelo liquida energía activa. De las 72 magnitudes que entregan
los equipos toma 2, y el Capítulo~\ref{sec:dato} explica por qué bastan
para simular un intercambio de energía. Queda por
responder algo distinto: cuánto pesa lo que se deja fuera, y si esa
omisión compromete la comparación entre mecanismos.

Este anexo responde a esa pregunta para la energía reactiva, que es la
magnitud omitida que sí tiene precio. La conclusión, adelantada, es que el
cargo es apreciable y el ordenamiento sobrevive.

\subsection{La energía reactiva y su cargo}
\label{sub:anexo-reactiva}

Aunque la energía reactiva no entrega trabajo, se paga. La regulación
colombiana la factura cuando el consumo reactivo inductivo
supera la mitad de la energía activa entregada en el mismo periodo
horario, y cobra el exceso como si fuera energía activa dentro de los
cargos por uso de las redes. El umbral equivale a
$\tan\varphi \le 0{,}5$, es decir, a un factor de potencia no menor que
\num{0,894}.

Medido sobre el horizonte, el exceso aparece en el \pct{15,0} de las horas
con consumo del circuito principal y en el \pct{65,7} de las del
secundario. El criterio de medida es el que fija la norma, que compara la
energía de cada hora y no los instantes sueltos. La distinción pesa menos
de lo que parece, porque contar las muestras de dos minutos deja el
resultado a menos de un punto porcentual, pero conviene declararla: son
dos preguntas distintas, la de cuánta energía incumple y la de cuántos
instantes lo hacen. La Tabla~\ref{tab:dato-reactiva} recoge las dos, por
institución y frontera de medición.

\begin{table}[H]
    \centering
    \small
    \caption{El exceso de energía reactiva sobre el umbral de la norma,
    por institución y frontera de medición, acumulado sobre el horizonte.}
    \label{tab:dato-reactiva}
    \begin{tabular}{@{}l r r r r r@{}}
    \toprule
    \textbf{Institución} & \textbf{Horas con} & \textbf{Horas en} &
    \textbf{Razón} & \textbf{Exceso} & \textbf{Cargo} \\
     & \textbf{consumo} & \textbf{exceso} & \textbf{mediana} &
    \textbf{(kvarh)} & \textbf{(COP)} \\
    \midrule
    \multicolumn{6}{@{}l}{\textit{M1 · circuito principal o de inyección}} \\
    Udenar  & \num{4491} & \num{284}  & \num{0,26} & \num{2217} & \num{470155} \\
    Mariana & \num{5737} & \num{3595} & \num{0,53} & \num{5876} & \num{1246650} \\
    UCC     & \num{5901} & \num{291}  & \num{0,27} & \num{661}  & \num{140271} \\
    HUDN    & \num{6031} & \num{1}    & \num{0,32} & \num{0}    & \num{2} \\
    CESMAG  & \num{6026} & \num{48}   & \num{0,05} & \num{37}   & \num{7938} \\
    \addlinespace
    \textbf{Total} & \num{28186} & \num{4219} & &
    \textbf{\num{8792}} & \textbf{\num{1865016}} \\
    \addlinespace
    \multicolumn{6}{@{}l}{\textit{M3 · circuito secundario}} \\
    Udenar  & \num{2354} & \num{0}    & \num{0,07} & \num{2}    & \num{526} \\
    Mariana & \num{187}  & \num{9}    & \num{0,13} & \num{105}  & \num{22228} \\
    UCC     & \num{6018} & \num{2662} & \num{0,41} & \num{461}  & \num{97858} \\
    HUDN    & \num{6031} & \num{6017} & \num{0,64} & \num{1606} & \num{340940} \\
    CESMAG  & \num{6024} & \num{4862} & \num{1,00} & \num{3531} & \num{748593} \\
    \addlinespace
    \textbf{Total} & \num{20614} & \num{13550} & &
    \textbf{\num{5705}} & \textbf{\num{1210145}} \\
    \bottomrule
    \end{tabular}
\end{table}

\notafig{La razón mediana y la fracción de horas en exceso no ordenan
igual a las instituciones, y por eso la tabla trae las dos. El Hospital
mantiene en el circuito principal una razón alta, \num{0,32}, y aun así no
rebasa el umbral ni una sola hora, mientras que Udenar lo rebasa en pocas
y acumula \uni{2234}{kvarh}, porque cuando lo rebasa lo hace por mucho. En
el circuito secundario, Mariana registra consumo en solo 187 horas y su
fila no admite la misma lectura que las demás.}

Que dos instituciones tengan tan pocas horas con consumo tiene tres causas
distintas, y conviene separarlas porque solo una es un defecto del dato. En
Udenar, sobre el circuito principal, el \pct{25,1} de las horas registra
potencia \emph{negativa}, hasta \uni{-33,6}{kW}. Ahí no falta consumo: hay
exportación, porque ese medidor es de los que descuentan la generación
propia de la lectura, y el Capítulo~\ref{sec:prep} se ocupa de
reconstruirlo. En su circuito secundario no hay ni una hora negativa, pero
el \pct{61} se queda entre 0 y \uni{0,5}{kW}, con una mediana de
\uni{0,38}{kW}: es sencillamente un ramal pequeño. El caso de Mariana en el
circuito secundario es de otro orden, con el \pct{90,4} de las horas por
debajo de \uni{0,5}{kW}, una mediana de \uni{0,02}{kW} y un máximo de
\uni{2,52}{kW}, que es lo que cabe esperar del medidor que el inventario
declara inexistente como segunda frontera.

La distinción importa para leer la tabla, porque sus dos columnas de horas
y las dos últimas no cuentan la misma población. La fracción de horas en
exceso se calcula solo sobre las horas con consumo apreciable, mientras que
la energía y el cargo recogen todas las horas. La norma lo exige así: cobra
el reactivo cuando supera la mitad de la activa y también cuando se
registra sin consumo activo alguno. En el circuito principal eso no es
menor, porque \uni{3473}{kvarh} de los \uni{8792}{kvarh}, es decir, el
\pct{39,5} del exceso facturable, nacen en horas que el recuento porcentual
deja fuera, y son en su mayoría las de exportación de Udenar.

\begin{detallebox}
Conviene precisar qué series son estas, porque el
Capítulo~\ref{sec:prep} somete el dato a 10 transformaciones antes de
entregarlo al modelo y cabe preguntar si lo que aquí se mide es
provisional.

No lo es, por dos razones distintas. La primera es que \textbf{no existe
una versión procesada de la energía reactiva}: el pipeline lee una sola
columna de cada medidor, la potencia activa total, de modo que la reactiva
nunca entra en él y la única serie disponible es la que el equipo
registró. La segunda es que, para la pregunta regulatoria, la lectura
cruda \emph{es} la correcta. La norma compara el reactivo contra la
energía activa entregada al usuario, y lo que se entrega al usuario es lo
que cruza el medidor, no lo que el edificio consume por dentro. La
reconstrucción del Capítulo~\ref{sec:prep} recupera esto último, que es lo
que el mercado necesita, pero no es lo que factura el comercializador.

La distancia entre las dos lecturas se puede medir, y resulta ser la
medida de lo que cuesta generar. Contra la demanda reconstruida, las horas
en exceso de Udenar bajan del \pct{6,3} al \pct{3,0}, las de Mariana del
\pct{62,7} al \pct{49,1} y las de la UCC del \pct{4,9} al \pct{0,2}, todas
en el circuito principal. El Hospital y la Universidad CESMAG no se mueven,
porque sus medidores no descuentan generación y no hay nada que
reconstruir. Es decir, las tres instituciones cuyo medidor resta la
generación propia incumplen más de lo que incumplirían por su consumo
real, desde un \pct{30} más en la Universidad Mariana hasta 22 veces más
en la UCC. El fotovoltaico no aumenta su demanda
de reactiva; reduce la activa contra la que se compara.
\end{detallebox}

\subsection{La forma del incumplimiento}
\label{sub:anexo-forma}

La Figura~\ref{fig:dato-reactiva-series} traza las 10 series
completas, una por institución y frontera, con la línea del umbral sobre
cada una, de modo que el
recorrido pueda verse entero en lugar de resumido.

\begin{figure}[H]
 \centering
 \includegraphics[width=\textwidth]{figuras/f2_05_reactiva_series.png}
 \caption[La razón entre reactiva y activa, hora a hora]{Razón entre
 energía reactiva y energía activa en cada hora con consumo, por
 institución y frontera. La mancha tenue son las horas sueltas y la línea
 gruesa la mediana de cada día. El área coloreada sobre la línea
 discontinua marca los tramos en que la norma cobra.}
 \label{fig:dato-reactiva-series}
\end{figure}

\notafig{La escala se recorta en \num{2,0}, lo que deja fuera el
\pct{0,7} de las horas. Las diez series tienen huecos, entre 97 y 126
horas de las \num{6144} del horizonte, con rachas de hasta dos días; el
Capítulo~\ref{sec:prep} explica cómo se rellenan.}

Lo que la serie muestra y una cifra de resumen no puede es la forma del
incumplimiento, que resulta ser distinta en cada caso. Mariana, en el
circuito principal, se instala sobre el umbral desde el primer día y no lo
abandona en ocho meses. La UCC oscila a diario y cruza el umbral en cada
ciclo, de modo que su \pct{44,1} en el circuito secundario no describe un
periodo malo sino la mitad de todos los días. El Hospital, en esa misma
frontera, dibuja una banda estrecha y constante justo por encima de la
línea, que es la forma que corresponde a una carga permanente.

La Universidad CESMAG merece una lectura aparte. En el circuito
secundario, su razón mediana diaria
vale \num{0,77} en abril y \num{0,80} en mayo, sube a alrededor de \num{1,0} en
junio y ya no baja de ahí en lo que queda del horizonte. Las horas en
exceso pasan del \pct{68} de mayo al \pct{98} de junio. Este trabajo no
puede determinar qué cambió en la instalación, y lo deja anotado porque un
sistema que empeora a mitad de horizonte y se queda así interesa a quien
lo opere.

\begin{figure}[H]
 \centering
 \includegraphics[width=0.95\textwidth]{figuras/f2_06_reactiva_estructura.png}
 \caption[El exceso según la hora del día]{Razón mediana entre energía
 reactiva y activa según la hora del día. La línea discontinua marca el
 umbral de la norma.}
 \label{fig:dato-reactiva-tiempo}
\end{figure}

\notafig{La serie de Mariana en el circuito secundario va punteada porque
allí solo el \pct{3} de las horas registra consumo.}

La Figura~\ref{fig:dato-reactiva-tiempo} añade el ciclo del día, que la
serie completa no deja leer. Mariana, en el circuito principal, alcanza su peor momento al mediodía,
con \num{0,95}. Su potencia activa mediana cae de \uni{9,1}{kW} a las 9
hasta \uni{6,2}{kW} a las 13, mientras que la reactiva baja solo de
\uni{5,8}{kvar} a \uni{5,3}{kvar}, de modo que la razón empeora sin que la
instalación consuma más reactiva. Es lo que cabe esperar cuando la
generación propia descuenta consumo activo y deja intacta la demanda de
reactiva. El Hospital y la Universidad CESMAG, en cambio, mantienen su
nivel a cualquier hora, lo que apunta a equipos que consumen reactiva de
forma continua.

\subsection{Cuánto cuesta, y si cambia el resultado}
\label{sub:anexo-cargo}

\begin{trampabox}
Puesto en dinero, el exceso suma \uni{1865016}{COP} en el circuito
principal y \uni{1210145}{COP} en el secundario, llevado a la tarifa de
transmisión y distribución, que promedia \uni{213}{COP/kWh} sobre el
horizonte.

La primera de esas dos cifras es la que obliga a declarar el asunto:
\textbf{el cargo que el modelo no representa supera el margen que separa a
los dos mecanismos mejor situados en esa frontera}, un margen de \uni{1426142}{COP}. No por ello cambia el ordenamiento. Es un cargo por
uso de redes que todos los mecanismos pagan sobre el mismo consumo, de
modo que en primera aproximación desplaza a todos por igual y las
diferencias sobreviven.

Queda un segundo orden que sí puede medirse. La norma compara el reactivo
contra la energía activa \emph{entregada al usuario}, con lo que un mecanismo que reduzca la energía que cada comprador toma de la red baja
también el umbral y convierte en facturable una parte mayor del reactivo.
El cálculo se hace institución por institución, que es como se aplica la norma, y bajo un supuesto que conviene declarar: cada kilovatio hora
comprado en el mercado desplaza uno de importación en el medidor del
comprador. Con él, el mercado eleva el cargo en \uni{97471}{COP} en el
circuito principal y en \uni{59782}{COP} en el secundario. Frente al
margen que lo separa del colectivo mensual, eso es el \pct{6,8} en el
primero, y frente a la brecha del segundo, el \pct{1,3}. Estrecha la
ventaja y no la invierte, de modo que el ordenamiento de los seis
mecanismos sobrevive intacto.
\end{trampabox}

\subsection{Por qué el trabajo se hace solo sobre energía activa}
\label{sub:anexo-porque}

\begin{detallebox}
De todo lo anterior se sigue la pregunta de fondo: si el reactivo cuesta
dinero y se comporta de manera tan desigual, por qué el trabajo entero se
hace sobre energía activa. Hay cuatro razones y conviene dejarlas dichas.

La primera es que \textbf{solo la activa transfiere energía}. Lo que el
mercado intercambia son kilovatios hora, y el reactivo no mueve energía
neta a lo largo de un ciclo, de modo que no hay nada que vender.

La segunda es que \textbf{los regímenes que se comparan están escritos en
energía activa}. La autogeneración, los excedentes y su reparto se miden
en kilovatios hora en las tres resoluciones que el Capítulo~\ref{sec:esc}
implementa. Comparar el mercado contra ellas sobre otra magnitud sería
comparar objetos distintos, que es justo lo que este documento evita.

La tercera es que \textbf{repartir reactivo exigiría representar la red}.
La compensación de reactivo es local, y afirmar que fluye de una
institución a otra necesita impedancias y un cálculo de flujos de
potencia. Este modelo no representa la red, y el Capítulo~\ref{sec:dato} ya
declaró que tampoco se afirma proximidad eléctrica entre los
participantes.

La cuarta es que \textbf{el modelo base está formulado sobre potencia
activa}. Incorporar el reactivo no sería añadir una columna sino reescribir
el juego, sus funciones de utilidad y su dinámica, y el trabajo perdería la
fidelidad al modelo que valida.

Lo que sí corresponde, y es lo que hace este anexo, es declarar el
cargo que queda fuera y medir si su omisión compromete la comparación. No
la compromete.
\end{detallebox}

\subsection{Una trampa de agregación}
\label{sub:anexo-agregacion}

\begin{trampabox}
Todas las cifras de este anexo están calculadas institución por
institución, y no sobre la suma de la comunidad. La diferencia no es de
matiz. Sumando primero la energía activa y la reactiva de las cinco, y
comparándolas después contra el umbral, el cargo del circuito principal
baja de \uni{1865016}{COP} a \uni{302108}{COP}, casi la sexta parte.

La razón es que la suma compensa lo que la norma no compensa. El exceso de
reactivo de una institución queda tapado por el consumo activo de otra, y
aparece una comunidad que cumple formada por miembros que no cumplen. La
resolución liquida a cada usuario contra su propia energía activa, de modo
que la única lectura válida es la desagregada.

Conviene tenerlo presente porque el error es fácil de cometer y difícil de
ver: la cifra agregada es más pequeña, más cómoda y perfectamente
plausible.
\end{trampabox}
